%% appE-manifest --- Manifest
%%
%% Generowany, nigdy nie pisany. Trzy rzeczy: manifest diagramów, manifest
%% ćwiczeń i rejestry zaległej pracy, drukowane zamiast tylko liczone przez
%% build, bo czytelnik z PDF-em nie uruchomi `make debt`.
%%
%% KAŻDA LICZBA NA TYCH STRONACH JEST OBLICZONA. code/measure/ledgers.py
%% zlicza je z drzewa i zapisuje figures/values/ledgers.tex, więc `make
%% verify` przerywa build, gdy rejestr się zmienia, a dodatek nie.

\chapter{Manifest}
\label{app:manifest}

\noindent
\emph{Co ta książka jest ci jeszcze winna?}

Każda liczba poniżej jest zliczona ze źródeł przez skrypt i trafia na stronę
mechanicznie, dokładnie tak jak każda zmierzona liczba w tej książce. To nie
porządek dla porządku: dodatek z ręcznie wpisanymi liczbami dezaktualizuje się
w tygodnie i nic tego nie widzi.

\section{Rejestr}
\label{sec:E-ledger}

\begin{center}
\begin{tabular}{@{}lr@{}}
\toprule
Rozdziały w manifeście & \val{ledger.chapters.total} \\
\quad napisane & \val{ledger.chapters.written} \\
\quad nadal zalążki, drukujące czerwoną ramkę \emph{\lblNotWritten} & \val{ledger.chapters.stubs} \\
Dodatki w manifeście & \val{ledger.appendices.total} \\
\quad nadal zalążki & \val{ledger.appendices.stubs} \\
Pliki z listingami w katalogu \code{code/} & \val{ledger.listings} \\
Ćwiczenia, każde z plikiem startowym, rozwiązaniem i testem & \val{ledger.exercises} \\
Listingi z ramką \emph{\lblVerify} & \val{ledger.verifybox} \\
Diagramy w tym wydaniu & \val{ledger.diagrams} \\
\bottomrule
\end{tabular}
\end{center}

\noindent
Liczby siedzą w tabeli, a nie w zdaniach, celowo. Zdanie musi się zgadzać ze
swoją liczbą \dash{} jeden plik, dwa pliki, pięć plików \dash{} a liczba, która
przychodzi ze skryptu, nie powie zdaniu, którą formę przyjąć; pierwszy szkic
tej strony mówił \enquote{1 plików}.

\section{Rozdziały}
\label{sec:E-chapters}

Rozdział, który nie jest napisany, drukuje czerwoną ramkę
\emph{\lblNotWritten} w miejscu tekstu, żeby liczba stron nigdy nie schlebiała
szkicowi. Ramka niesie kontrakt rozdziału \dash{} co ma uzasadnić i czego nie
ma powtarzać \dash{} więc jest umową, a nie zapchajdziurą.

\section{Listingi i ćwiczenia}
\label{sec:E-listings}

Każdy plik z listingiem w katalogu \code{code/} i każde ćwiczenie jest
uruchamiane przez build repozytorium przy każdej zmianie, na wersjach ze
strony tytułowej. Listing, który nosi ramkę \emph{\lblVerify}, przez ten bieg
nie przeszedł; jego liczba jest wydrukowana powyżej po to, żeby nie dało się
jej po cichu nieść dalej.

\listofexercises

\section{Diagramy}
\label{sec:E-diagrams}

Każdy diagram jest narysowany osobno w każdym języku, bo diagram z
angielskimi etykietami w polskim wydaniu to ten rodzaj półprzekładu, przez
który dwujęzyczna książka wygląda na zrobioną maszynowo. Źródłem prawdy jest
plik tekstowy, wersjonowany, więc zmiana jest widoczna jako diff. Listę
poniżej pisze makro figury w miejscu, w którym figura jest umieszczona, więc
diagram nie może być w książce i brakować na liście.

\listofdiagrams

\section{Pomiary}
\label{sec:E-measurements}

W rozdziałach wyspecyfikowano osiem pomiarów, wszystkie darmowe i kończące
się na laptopie. Każdy rozdział, który ma swój pomiar, nazywa go w swoim
kontrakcie, a zdanie, które pomiar miałby wspierać, jest oznaczone jako osąd,
dopóki skrypt nie został uruchomiony. Nie drukuje się tu sumy, ile z nich
uruchomiono, celowo: taka liczba jest twierdzeniem o książce, którego nic nie
wyprowadza z niczego, i dezaktualizuje się po cichu.
